\section{FP3 --- Écoute et validation de l'enregistrement}
\label{sec:fp3}

Fonction couverte par l'exigence \etref{4} : \emph{une lecture sur Audacity doit pouvoir restituer de manière claire l'enregistrement du mot « Électronique »}.
\subsection{Objectif}

Cette fonction sert de \emph{contrôle de bout en bout} des étages FP1 et FP2 : si l'enregistrement reproduit fidèlement la voix humaine à l'écoute, cela atteste que la numérisation et le filtrage fonctionnent correctement. C'est une fonction de validation, non un bloc fonctionnel du produit final.

\subsection{Conception}

\subsubsection{Protocole de transfert}

\TODO{Décrire :
\begin{itemize}
    \item Déclenchement par appui bouton $\rightarrow$ capture d'\SI{1}{\second} = 8000 échantillons à \SI{8}{\kilo\hertz}
    \item Format des échantillons : int16 little-endian (2 octets / échantillon)
    \item Baudrate série : \num{250000} bauds
    \item Header de synchronisation : \TODO{ex. \texttt{0xAA55} ou caractère délimiteur}
\end{itemize}}

\subsubsection{Script Python côté PC}

\TODO{Commit du script \texttt{scripts/serial\_to\_wav.py} qui :
\begin{enumerate}
    \item Ouvre le port série
    \item Attend le header
    \item Lit 8000 int16
    \item Écrit un fichier \texttt{.wav} (16 bits, mono, \SI{8}{\kilo\hertz}) via \texttt{scipy.io.wavfile} ou \texttt{wave}
\end{enumerate}
Inclure l'extrait clé du script dans le rapport.}

\subsection{Validation}

\subsubsection{Enregistrement du mot « Électronique »}

\TODO{Figure \texttt{fp3-audacity-electronique.png} : capture Audacity montrant :
\begin{itemize}
    \item La forme d'onde temporelle du mot (visualisation des 4 syllabes)
    \item Le spectrogramme associé (bande 0--\SI{4}{\kilo\hertz})
    \item La fréquence d'échantillonnage indiquée (\SI{8}{\kilo\hertz})
\end{itemize}}

\subsubsection{Validation auditive}

\TODO{Décrire l'écoute : intelligibilité, absence de cliquetis, absence d'artefacts de sous-échantillonnage. Stocker le fichier \texttt{.wav} dans \texttt{docs/audio-samples/electronique.wav}.}

\subsubsection{Comparaison spectrale avant/après filtre}

\TODO{Optionnel mais valorisant : comparer le spectre du signal brut (si on enregistre directement la sortie ADC à \SI{32}{\kilo\hertz}) vs le signal filtré \& sous-échantillonné à \SI{8}{\kilo\hertz}. Montre visuellement l'effet du filtre.}

\subsection{Bilan FP3}

\TODO{Synthèse : la chaîne FP1 + FP2 est validée. Le mot est audible et intelligible dans Audacity.}

\newpage
